\chapter*{LEMBAR PENGESAHAN}
\addcontentsline{toc}{chapter}{LEMBAR PENGESAHAN}


\begin{center}
    \vspace{2cm}
  {\large\bfseries Perbandingan Model \emph{Machine Learning} dan Metode \emph{Explainable AI} untuk Prediksi \emph{Cumulative Abnormal Return} Emiten IDX80 Berdasarkan Jenis Peristiwa pada Berita Berbahasa Indonesia}\\
     \vspace{2cm}

  {\Large \textbf{Proposal Tugas Akhir}}\\


  \vspace{1.5cm}
    
    
  {\large Oleh}\\[0.3cm]
    \textbf{
    {\large SEBASTIAN ALBERN NUGROHO}\\
    {\large 18223074}
  }\\
    
  \vspace{0.5cm}
 
  {\large Program Studi Sistem dan Teknologi Informasi}\\
  {\large Sekolah Teknik Elektro dan Informatika}\\
  {\large Institut Teknologi Bandung}\\

  \vspace{1.5cm}

  Proposal Tugas Akhir ini telah disetujui dan disahkan\\
  di Bandung, pada tanggal \today\\[1cm]

% ==========================================
% Versi 1 pembimbing (default)
% ==========================================
	Pembimbing  \\[3cm]
	Dr. Nur Ulfa Maulidevi, S.T., M.Sc.   \\[0.2cm]
	NIP. 197603092008012010 
% ==========================================

\end{center}

\vspace{1cm}
\noindent

% ==========================================
% Jika ada 2 pembimbing TA, uncomment dan edit 
% tabular di bawah ini. Kemudian, comment out atau hapus
% bagian versi 1 pembimbing di atas.
% ==========================================

%\begin{tabular}{p{1cm}p{7cm}p{7cm}}
%   & Pembimbing 1 & Pembimbing 2 \\[3cm]
%   & Dr. Ir. Nama Pembimbing, M.T. & Dr. Nama Pembimbing, M.Sc. \\[0.2cm]
%   &  NIP. 123456789 & NIP. 987654321
%\end{tabular}
